\documentclass[a4paper]{article}
\usepackage{amsfonts}
\usepackage{amsmath}
\usepackage{amssymb}
\usepackage{graphicx}     

\begin{document}
  
\noindent \large Auteur: Michael Zielinski \\
\noindent \large Studierichting: Informatica\\
\noindent \large Oefeningenreeks 6F nr. 10\\

\medskip

\normalsize

$ f_n:\mathbb{R}\rightarrow\mathbb{R}:x\mapsto e^{\tfrac{x}{n}} $ \\

We bepalen eerst de puntsgewijze limiet. \\

Voor een vaste $x\in\mathbb{R}$ geldt: \\

$ \lim\limits_{n\rightarrow\infty} f_n(x)
= \lim\limits_{n\rightarrow\infty} e^{\tfrac{x}{n}} $ \\

Omdat: \\

$ \lim\limits_{n\rightarrow\infty}\dfrac{x}{n}=0 $ \\

volgt: \\

$ \lim\limits_{n\rightarrow\infty} e^{\tfrac{x}{n}}=e^0=1 $ \\

Dus: \\

$ f:\mathbb{R}\rightarrow\mathbb{R}:x\mapsto 1 $ \\

$ f_n \stackrel{p}{\rightarrow} f $ \\

Nu onderzoeken we uniforme convergentie. \\

$ \left|\left|f_n-f\right|\right|_{\infty}
= \sup\limits_{x\in\mathbb{R}}\left|e^{\tfrac{x}{n}}-1\right| $ \\

Voor $x\rightarrow+\infty$ geldt: \\

$ e^{\tfrac{x}{n}}\rightarrow+\infty $ \\

Dus: \\

$ \left|e^{\tfrac{x}{n}}-1\right|\rightarrow+\infty $ \\

Hieruit volgt: \\

$ \sup\limits_{x\in\mathbb{R}}\left|e^{\tfrac{x}{n}}-1\right|=+\infty $ \\

Dus: \\

$ \left|\left|f_n-f\right|\right|_{\infty}=+\infty $ \\

Daarom geldt: \\

$ \lim\limits_{n\rightarrow\infty}\left|\left|f_n-f\right|\right|_{\infty}\neq 0 $ \\

Dus $f_n$ convergeert niet uniform naar $f$. \\

Conclusie: \\

$ f_n \stackrel{p}{\rightarrow} f $ \\

maar \\

$ f_n \stackrel{u}{\not\rightarrow} f $ \\

\begin{thebibliography}{999}
%\bibitem{a} Referentie
\end{thebibliography}

\end{document}